\documentclass[12pt]{article}
\usepackage{xeCJK}  

\usepackage{charter}
\usepackage{fullpage}
\usepackage[colorlinks=false]{hyperref}
\usepackage{ifthen}
\usepackage{comment}
\usepackage[title,titletoc]{appendix}
\usepackage{pagecolor}
\usepackage{amsmath}
\usepackage{amsfonts}
%\usepackage[normalem]{ulem}
\usepackage{siunitx}
\usepackage{amsthm}
\sisetup{per=slash, load=abbr}

\usepackage{pgfplots}
\usetikzlibrary{positioning}
\usetikzlibrary{fit}
\usetikzlibrary{snakes}
\usetikzlibrary{shapes.geometric}
\usetikzlibrary{patterns}
\usetikzlibrary{shapes,arrows,chains}
\usepgfplotslibrary{patchplots,colormaps}
\usetikzlibrary{calc}
\usetikzlibrary{positioning, fit}
\usetikzlibrary{backgrounds}
\usetikzlibrary{intersections}

\newcommand{\whitepaper}[1]{\begin{center}\fbox{\parbox{0.95\textwidth}{{\footnotesize
#1}}}\end{center}}

\newcommand{\pcolor}{white!25}

\usepackage{setspace}
\usepackage{algorithm2e}
\bibliographystyle{ieeetr}

\usepackage{geometry}
\geometry{left=3cm,right=3cm,top=1.6cm,bottom=3cm,headheight=0pt,headsep=1.5em}
\usepackage{fancyhdr}
\pagestyle{fancy}
\chead{\includegraphics[scale=0.1]{../../common/Nebulas.png}}  %在此处插入logo.pdf图片 图片靠左
\lhead{} % 页眉中间位置内容
\rhead{}
%\setlength{\topskip}{1em}

\usepackage{indentfirst}


\setCJKmainfont[BoldFont = STSongti-SC-Bold]{STSongti-SC-Regular}
\setCJKfamilyfont{hei}{SIL-Hei-Med-Jian}		%宋体
\setCJKfamilyfont{song}{SimSun}		%宋体
\setCJKfamilyfont{kai}{Kaiti}		%楷体
\setCJKfamilyfont{fang}{song}	%仿宋
\setCJKfamilyfont{li}{song}			%隶书
\setCJKfamilyfont{you}{Yuanti}		%幼圆

\newcommand{\song}{\CJKfamily{song}}	%宋体
\newcommand{\hei}{\CJKfamily{hei}}	%黑体
\newcommand{\kai}{\CJKfamily{kai}}	%楷体
\newcommand{\fs}{\CJKfamily{fang}}	%仿宋
\newcommand{\li}{\CJKfamily{li}}		%隶书
\newcommand{\you}{\CJKfamily{you}}	%幼圆
\newcommand{\reffig}[1]{图\ref{#1}}
\newcommand{\refsec}[1]{\S \ref{#1}}

\onehalfspacing   % ----------设置1.5倍行距（可能有意义，待调整）

%\parindent=20pt  % -------------------首行缩进大小，英文分段就直接0pt了吧。
\setlength{\parindent}{2.1em}
\setlength{\parskip}{0.3\baselineskip}
\newcommand{\nrcore}{Core Nebulas Rank}
\newcommand{\nrext}{Extended Nebulas Ranks}
\newcommand{\dom}{{\; \texttt{dom}\;}}
\newcommand{\xpcomment}[1]{{\color{blue} #1}}

\setCJKsansfont[BoldFont = STHeitiSC-Medium]{STHeitiSC-Light}


\newtheorem{property}{特征}
%\addbibresource{reference.bib}

\begin{document}
\pagestyle{empty}
\renewcommand{\contentsname}{目录}
\renewcommand{\abstractname}{摘要}
\renewcommand{\refname}{参考文献}
%\renewcommand{\nomname}{术语表（按首字母排序）}
\renewcommand{\figurename}{图}
\renewcommand{\tablename}{表}
\renewcommand{\baselinestretch}{1.5}
\renewcommand{\appendixname}{附录}
\renewcommand{\proofname}{证明}

\pagecolor{\pcolor}

\begin{titlepage}
  \begin{center}
    \vspace*{5.5cm}
    \includegraphics[scale=0.4]{../../common/Nebulas.png}
    \vspace{0.5cm}


    \textbf{\huge{星云临时技术委员会章程}}

    \vspace{0.5cm}
    
    \vfill
    2019年6月 \\
    版本号:0.1
    \textbf{}
  \end{center}

\end{titlepage}
\setcounter{page}{0}
%\thispagestyle{empty}
\tableofcontents
\newpage
\setcounter{page}{1}
\pagestyle{fancy}	
\vspace*{0.01cm}
\section{介绍}
\subsection{星云技术委员会}

星云技术委员会成立于2018年9月。成立至今，星云技术委员会秉持着开放、共享、透明的精神，致力于推动星云链技术研发逐步去中心化、社区化。

自星云理事会成立后，原由星云团队核心成员构成的星云技术委员会将完成历史使命，转型为社区化团队。星云技术委员会受星云理事会委托，负责星云项目制生产力组织及质量核查，为社区提供技术指引和支持。

\subsubsection{人员组成}

星云技术委员会人数不限。

\subsubsection{权力}

\begin{enumerate}
	\item 向理事会发起提案重审的权力。
	\item 享有星云技术专家团队的荣誉。
\end{enumerate}

\textbf{义务}

\begin{enumerate}
	\item 社区化提案质量监督，出具相关测试报告和技术评级报告。
\end{enumerate}

\subsubsection{任期}

技术委员会委员任期1年，可连任。

\subsubsection{入选方法}

技术委员会采用自荐和社区推荐相结合的方法选择，对社区公开述职。具体将由星云理事会组织进行。

\subsubsection{收益}

\textbf{总收入}

\begin{itemize}
	\item 委托佣金（按月发放）
	\item 项目审核和监督的咨询费用
\end{itemize}

\vspace{2em}

\textbf{财务要求}

为保证经济体利益的一致性，和星云技术委员会政策的延续性，星云技术委员会委员正式入职时需要抵押25,000 NAS。卸任3个月后解锁返还。

\end{document}
